\begin{document}
\maketitle
\begin{abstract}
An image pair can now produce a dense 3D Gaussian reconstruction in one
network pass. Building a lasting map from these predictions is harder:
surfaces overlap, and loop closures move earlier camera poses.
We present Splatt3R-SLAM, which connects Splatt3R prediction to
MASt3R-SLAM through shared keyframe anchors. Each anchor carries both a local
Gaussian map and the camera poses used to refine it. Pose corrections
therefore move the map and its observations together. A separate refiner
improves appearance while the geometric tracker estimates motion.
On eight Replica scenes, the system reaches 23.02\,dB mean PSNR after
300\,s refinement, compared with 19.48\,dB for Photo-SLAM under common
non-keyframe reconstruction evaluation. On TUM desk, PSNR is close to MonoGS,
which has lower LPIPS. Component studies isolate the effects of refinement,
head adaptation, and opacity attenuation. Dense prediction provides a useful
map initialisation, while its millions of primitives make compact online
mapping the main remaining challenge.
\end{abstract}
\ifralreview\else
\begin{IEEEkeywords}
SLAM, mapping, visual learning, 3D Gaussian splatting.
\end{IEEEkeywords}
\fi

\input{ral/introduction}
\input{ral/method}
\input{ral/experiments}
\input{ral/discussion}
\bibliographystyle{IEEEtran}
\bibliography{main}
\end{document}
